\documentclass[11pt]{article}

\usepackage[letterpaper,margin=1in]{geometry}
\usepackage{amsmath,amssymb}
\usepackage{booktabs}
\usepackage{array}
\usepackage{xcolor}
\usepackage{enumitem}
\usepackage{titlesec}
\usepackage{hyperref}
\usepackage{parskip}
\usepackage{microtype}
\usepackage{longtable}
\sloppy

\definecolor{notegray}{HTML}{666666}
\definecolor{headshade}{HTML}{D9E2F3}

\newcommand{\note}[1]{{\itshape\color{notegray}#1\par}}
\newcommand{\code}[1]{\texttt{#1}}

\titleformat{\section}{\normalfont\Large\bfseries}{\thesection}{1em}{}
\titleformat{\subsection}{\normalfont\large\bfseries}{\thesubsection}{1em}{}

\setlist[itemize]{leftmargin=1.4em, itemsep=0.35em, topsep=0.3em}

\hypersetup{colorlinks=true, linkcolor=blue, urlcolor=blue, citecolor=blue}

\title{\textbf{DL4DR Compound Generator --- RL Fine-Tuning}\\[4pt]
\Large Complete Investigation: Reward Hacking, Mode Collapse, and\\Specification Gaming Across v1--v5}
\author{}
\date{\itshape\color{notegray} Consolidated report --- 2026-07-31}

\begin{document}
\maketitle
\vspace{-1.5em}

\begin{abstract}
\noindent
This report consolidates a five-run investigation (v1--v5) into REINFORCE fine-tuning of a conditional SMILES VAE compound generator against a frozen drug-response predictor (DL4DR). The starting question was whether the policy learns to genuinely improve predicted drug potency or instead learns to \emph{exploit} the reward signal. No classic reward hacking against the predictor itself was ever found (the frozen model was never fooled into favoring non-drug-like molecules). What was found, characterized, and eventually fixed was a subtler and arguably more instructive failure mode: \emph{specification gaming against an auxiliary diversity reward}, appearing in two successive forms (v1--v3's mode collapse, then v4's scaffold-repetition loophole), each diagnosed mathematically and each addressed by a different mechanism (v4's exact-match repetition penalty, then v5's Tanimoto-similarity-clustered version). Part~I covers the setup and v1--v4; Part~II covers v5's fix and final verdict.
\end{abstract}

\tableofcontents
\newpage

\part{Setup and v1--v4}

\section{Objective}

This project fine-tunes a conditional SMILES VAE compound generator against a frozen DL4DR drug-response predictor (\code{best\_random.pt}) as a reward model, using REINFORCE. The specific question under investigation: does the policy learn to genuinely generate better drug candidates (lower predicted IC50), or does it learn to exploit (``hack'') the reward signal in some way that does not reflect real improvement?

\section{Method Summary}

\begin{itemize}
  \item Reward $= -\text{predicted}\ \ln(\text{IC50})$ from the frozen predictor, plus a fixed penalty for invalid/undecodable SMILES.
  \item Full-VAE variant (\code{new\_rl\_finetune.py}): both the encoder and decoder are trained. Each epoch, real seed SMILES for the current cell line are encoded to a posterior $z$ (instead of sampling $z$ from a fixed prior), so REINFORCE gradients reach the encoder as well as the decoder. A standard VAE ELBO term (reconstruction $+$ KL) is added on the same $z$ to keep the encoder's posterior well-behaved.
  \item Independent evaluation (\code{new\_rl\_eval\_independent.py}) scores generated molecules two ways at once: (a)~the frozen predictor's own reward-model score --- what the policy was optimized for; and (b)~reward-model-blind checks the predictor cannot see or influence --- RDKit validity, Lipinski Rule-of-Five druglikeness, novelty vs.\ the training set, nearest-neighbor Tanimoto similarity, and uniqueness. A widening gap between (a) climbing and (b) staying flat or falling is the reward-hacking signal this report is checking for.
\end{itemize}

\section{Metric Definitions}

These four independent metrics measure different, non-overlapping properties of the generated molecules:

\begin{center}
\begin{tabular}{@{}l p{0.72\textwidth}@{}}
\toprule
\textbf{Metric} & \textbf{What it measures} \\
\midrule
Validity & Can this single SMILES string be parsed into a chemically valid molecule (correct atom valence, closed rings)? A per-sample, syntactic check only. \\
Uniqueness & Of the valid molecules in one batch, what fraction are structurally distinct from each other? A within-batch diversity check --- independent of validity. \\
Novelty rate & Of the unique valid molecules, what fraction do NOT already appear in the training set? Checks for memorization/copying vs.\ genuine generation. \\
Lipinski pass rate & Standard Rule-of-Five druglikeness filter (molecular weight, LogP, H-bond donors/acceptors, rotatable bonds). The predictor was never trained on this rule, so it is a reward-model-blind sanity check. \\
Reward-model score & $-\text{predicted}\ \ln(\text{IC50})$ from the frozen predictor --- the actual quantity the policy was optimized to maximize. \\
\bottomrule
\end{tabular}
\end{center}

\textbf{Key point:} high validity does not imply high molecule quality or diversity. It only means the generated string is chemically parseable. A policy can reach $\sim$100\% validity by repeatedly emitting the same few known-safe molecules --- which is exactly what is observed below (Section~\ref{sec:results}).

\section{Runs Compared (v1--v4)}

\begin{center}
\begin{tabular}{@{}l l l l l@{}}
\toprule
 & \textbf{v1} & \textbf{v2} & \textbf{v3} & \textbf{v4} \\
\midrule
Epochs & 1000 & 2000 & 2000 & 2000 \\
\code{entropy\_coef} & 0.01 & 0.05 & 0.5 & 0.05 \\
\code{invalid\_penalty} & $-10.0$ & $-5.0$ & $-5.0$ & $-5.0$ \\
\code{lambda\_div} & --- & --- & --- & 0.1 (new) \\
Resumed from & pre-RL VAE & pre-RL VAE & pre-RL VAE & pre-RL VAE \\
\bottomrule
\end{tabular}
\end{center}

\note{All four runs resume from the same pre-RL warm-start VAE checkpoint (not from each other), so v1--v4 are a clean controlled comparison isolating \code{entropy\_coef}, \code{invalid\_penalty}, and \code{lambda\_div}. v4 is the first to add the Section~\ref{sec:penalty} batch-level repetition penalty on top of v2's hyperparameters (the best-performing combination through v3).}

\section{Results}
\label{sec:results}

\subsection{Training-log trend (v1 vs.\ v2, last $\sim$50 logged epochs)}

\begin{center}
\begin{tabular}{@{}l l l@{}}
\toprule
\textbf{Metric} & \textbf{v1 tail (epoch $\sim$950--1000)} & \textbf{v2 tail (epoch $\sim$1947--2000)} \\
\midrule
Entropy & 0.06--0.24 & 0.025--0.091 (lower) \\
Validity & 92--100\% & 93--100\% (similar) \\
\code{mean\_reward\_valid} & $-3.7 \sim -4.0$ & $-2.7 \sim -3.35$ (better) \\
\bottomrule
\end{tabular}
\end{center}

\subsection{Independent evaluation (500 samples/cell line, same 3 cell lines, both sampling modes)}
\label{sec:eval_v1v4}

\begin{center}
\small
\begin{tabular}{@{}l l l l l l@{}}
\toprule
\textbf{Mode} & \textbf{Metric} & \textbf{v1} & \textbf{v2} & \textbf{v3} & \textbf{v4} \\
\midrule
prior & Reward-model score & $-3.71\sim-4.31$ & $-3.47\sim-3.55$ & $-3.79\sim-4.03$ & $-3.66\sim-3.81$ \\
 & Validity & 48.8--57.2\% & 58.8--65.4\% & 54.4--57.0\% & 65.4--68.0\% \\
 & Uniqueness & 37.8--38.9\% & 24.2--27.2\% & 24.5--26.1\% & 51.2--56.9\% \\
posterior & Reward-model score & $-3.70\sim-3.90$ & $-3.14\sim-3.24$ (best) & $-3.63\sim-3.82$ & $-3.39\sim-3.61$ \\
 & Validity & 97.8--100\% & 97.0--98.4\% & 96.0--96.4\% & 96.0--97.4\% \\
 & Uniqueness & 6.3--7.8\% & 5.1--7.0\% & 5.6--6.9\% & \textbf{28.2--31.6\%} (best) \\
both & Novelty rate & 100\% & 100\% & 100\% & 100\% \\
both & Lipinski pass rate & 100\% & 100\% & 100\% & 99.3--100\% \\
both & mean\_NN\_Tanimoto & 0.26--0.30 & 0.27--0.30 & 0.27--0.32 & 0.27--0.30 \\
\bottomrule
\end{tabular}
\end{center}

\textbf{v4 is the first run with both good validity and good diversity at once:} posterior-mode uniqueness (28.2--31.6\%) is 4--5x higher than v1/v2/v3 (5.1--7.8\%), while validity (96.0--97.4\%) and reward ($-3.39\sim-3.61$) stay in the same competitive range as v2/v3. \code{entropy\_coef} tuning alone (v1$\to$v2$\to$v3) never achieved this; adding the batch-level repetition penalty (Section~\ref{sec:penalty}) on v2's base hyperparameters did.

\subsection{v3 training-log trajectory (entropy\_coef $=0.5$, applied from epoch 1)}
\label{sec:v3traj}

Unlike v2 (\code{entropy\_coef} raised mid/late training), v3 applies a 50x larger \code{entropy\_coef} ($0.5$ vs.\ v1's $0.01$) from the very first epoch --- a direct test of the hypothesis (Section~\ref{sec:8.5}) that early application should avoid the softmax-saturation problem.

\begin{center}
\begin{tabular}{@{}r r r r r@{}}
\toprule
\textbf{Epoch} & \textbf{Entropy} & \textbf{Validity} & \textbf{mean\_reward\_valid} & \textbf{$|$rl\_loss$|$} \\
\midrule
1 & 1.66 & 10.9\% & $-4.48$ & 993.5 \\
100 & 0.62 & 26.6\% & $-4.12$ & 11.9 \\
250 & 0.29 & 73.4\% & $-4.26$ & --- \\
500 & 0.12 & 82.8\% & $-4.05$ & 1.69 \\
750 & 0.08 & 85.9\% & $-3.00$ & --- \\
1000 & 0.10 & 89.1\% & $-3.89$ & 2.96 \\
1500 & 0.10 & 95.3\% & $-3.84$ & --- \\
2000 (tail avg, last 20) & 0.055 & 95.1\% & $-3.78$ & 2.27 \\
\bottomrule
\end{tabular}
\end{center}

\textbf{Result: the 50x larger, epoch-1 \code{entropy\_coef} did NOT prevent collapse.} Entropy was already down to 0.10--0.12 by epoch 500--1000 --- statistically indistinguishable from where v1 (\code{entropy\_coef}$=0.01$) ended up by \emph{its} epoch 1000 (0.06--0.24). By the tail, v3's entropy (avg $0.055$) is in the same collapsed range as v2's, and v3's reward (avg $-3.78$) is actually \emph{worse} than v2's ($-2.7\sim-3.35$): a 10x larger \code{entropy\_coef} than v2 bought no diversity benefit and cost some reward quality. This falsifies the simple form of the Section~\ref{sec:8.5} recommendation and is resolved with a revised mechanism in Section~\ref{sec:8.7}.

\subsection{v4 training-log trajectory (\code{lambda\_div}$=0.1$ added on top of v2's hyperparameters)}

\begin{center}
\begin{tabular}{@{}r r r r r@{}}
\toprule
\textbf{Epoch} & \textbf{Entropy} & \textbf{Validity} & \textbf{batch\_uniqueness} & \textbf{mean\_reward\_valid} \\
\midrule
1 & 1.663 & 10.9\% & 71.4\% & $-4.48$ \\
100 & 0.567 & 28.1\% & 83.3\% & $-4.13$ \\
250 & 0.180 & 68.8\% & 93.2\% & $-4.24$ \\
500 & 0.163 & 89.1\% & 82.5\% & $-4.12$ \\
750 & 0.182 & 76.6\% & 69.4\% & $-2.99$ \\
1000 & 0.206 & 90.6\% & 65.5\% & $-3.56$ \\
1500 & 0.146 & 98.4\% & 63.5\% & $-3.45$ \\
1750 & 0.158 & 95.3\% & 67.2\% & $-3.53$ \\
2000 (tail avg, last 25) & 0.134 & 95.9\% & 55.3\% & $-3.44$ \\
\bottomrule
\end{tabular}
\end{center}

\textbf{Entropy stayed just as low as v1/v2/v3 (tail avg 0.134) --- but \code{batch\_uniqueness} stayed high (tail avg 55.3\%) throughout training, instead of collapsing to single digits.} This is the clearest evidence yet that per-step entropy and batch-level output diversity are different axes (mechanism explained in Section~\ref{sec:scaffold}): the repetition penalty fixed diversity directly, without needing the decoder to become less confident at each token.

\subsection{v4 broader independent evaluation (15 cell lines)}

\begin{center}
\begin{tabular}{@{}l l l l@{}}
\toprule
\textbf{Mode} & \textbf{Metric} & \textbf{Range (15 cell lines)} & \textbf{Mean} \\
\midrule
prior & Reward-model score & $-3.50\sim-3.81$ & $\approx -3.71$ \\
 & Validity & 61.6--69.4\% & $\approx 65.7\%$ \\
 & Uniqueness & 47.7--56.9\% & $\approx 51.9\%$ \\
posterior & Reward-model score & $-3.30\sim-3.61$ & $\approx -3.45$ \\
 & Validity & 93.8--98.2\% & $\approx 96.1\%$ \\
 & Uniqueness & 24.6--31.8\% & $\approx 28.9\%$ \\
both & Novelty / Lipinski pass rate & 99.4--100\% / 97.8--100\% & $\approx 100\%$ / $\approx 99.6\%$ \\
\bottomrule
\end{tabular}
\end{center}

Confirms Section~\ref{sec:eval_v1v4} was not a 3-cell-line fluke: posterior uniqueness averages $\approx$28.9\% across 15 cell lines, consistent with the 3-cell-line number.

\subsection{Important caveat: scaffold-level repetition}
\label{sec:scaffold}

The top-repeated-molecule diagnostic reveals a pattern invisible in the aggregate uniqueness numbers above: in posterior mode, across ALL 15 cell lines evaluated, the top-2 most-repeated molecules are the same two structures every time:

\begin{center}
\begin{tabular}{@{}l l l@{}}
\toprule
\textbf{Canonical SMILES} & \textbf{Top-2 repeat in} & \textbf{Typical count (of 500)} \\
\midrule
\code{CCCOC(O)OC} & 15 of 15 cell lines & 35--59 \\
\code{COC(O)OCCO} & 15 of 15 cell lines & 28--55 \\
\bottomrule
\end{tabular}
\end{center}

Close relatives of the same template (\code{CCCOC(O)OCO}, \code{CCCOC(O)OCCO}, \code{OCCOC(O)OCCO}, \code{OCCOC(O)OCO} --- chain-length/branching variants of the same $[\text{alkyl}]$--O--CH(OH)--O--$[\text{alkyl}]$ scaffold) fill out most of the remaining top-5 slots in every single cell line's report.

\textbf{Root cause: the repetition penalty (Section~\ref{sec:penalty}) only matches EXACT canonical-SMILES duplicates.} \code{compute\_diversity\_penalty()} computes $\text{count}_i$ by exact string equality --- \code{CCCOC(O)OC} and \code{CCCCOC(O)OC} (one more carbon) are, correctly per the code, two different strings, so neither is penalized relative to the other. The policy found it can satisfy the \emph{letter} of the uniqueness reward by enumerating chain-length/branching variants of one simple, easy-to-decode scaffold, rather than by exploring genuinely different chemistry --- a subtler, second-order form of specification gaming.

\textbf{This is not a data-coverage problem:} training already visits all 51 cell lines via \code{epoch \% 51} ($2000$ epochs $\Rightarrow \approx 39$ visits/cell line), identically across v1--v4 --- the scaffold-reuse pattern recurring across every cell line is a property of the penalty's exact-match granularity, not of how many cell lines or epochs were used. See Part~II for the fix.

\section{Interpretation}

\begin{itemize}
  \item Reward genuinely improved in v2, in both sampling modes. This is real progress on the actual objective, not just a validity artifact --- unlike v1, where reward for valid molecules was essentially flat across training despite validity climbing from $\sim$10\% to $\sim$97\%.
  \item Entropy collapsed further in v2, not less, despite \code{entropy\_coef} increasing 5x --- explained rigorously in Section~\ref{sec:8.4}.
  \item Prior-mode uniqueness got worse in v2 (38\%$\to$25\%) even as prior-mode validity improved --- consistent with the whole policy narrowing onto a smaller, higher-reward set of outputs, not a mode-specific effect.
  \item Posterior-mode collapse persists essentially unchanged ($\sim$6--8\% uniqueness) across v1--v3: the encoder-seeded path generates a small, repeated set of molecules regardless of which real compound was encoded as the seed.
  \item No strict reward-hacking flag (predictor rewarding non-drug-like molecules) fired in any run: Lipinski pass rate stayed near 100\% and reward-model score stayed negative throughout.
\end{itemize}

\section{Verdict (v1--v4)}
\label{sec:verdict1}

\textbf{No classic reward hacking detected} (predictor exploited via non-drug-like molecules). Lipinski pass rate held at $\geq$99\% and reward-model score remained solidly negative across all four runs and both sampling modes --- the specific failure mode the evaluation script is built to catch did not occur.

\textbf{A related but distinct failure mode was found instead: specification gaming via validity, compounded by mode collapse.} Over 1000--2000 epochs, the policy learned to reliably avoid \code{invalid\_penalty} (validity 10\%$\to\sim$97--100\%) largely by narrowing its output distribution onto a small, repeated set of molecules, rather than by broadly exploring toward better drug candidates.

\textbf{Confirmed by a pre-RL baseline (Section~\ref{sec:baseline}):} the collapse is caused by RL fine-tuning, not inherent to the VAE --- pre-RL posterior-mode uniqueness is 45.8--53.9\%, dropping to 6.3--7.8\% after RL.

\textbf{v4 update:} the batch-level repetition penalty (Section~\ref{sec:penalty}) recovers most of the lost diversity (uniqueness $\to$28.2--31.6\%, 4--5x v1/v2/v3) at v2-level reward/validity --- but Section~\ref{sec:scaffold}'s diagnostic shows this is not yet clean: the same scaffold family dominates every cell line, because the penalty only catches exact string duplicates, not near-identical chain-length variants. v4 is real progress, not a solved problem.

\section{Theoretical Analysis: Role of Entropy in REINFORCE Training}
\label{sec:theory}

\subsection{Setup --- matching the codebase}

Let $\pi_\theta$ denote the decoder policy over token sequences $\tau = (a_1, \ldots, a_T)$, conditioned on cell-line embedding $z_C$ and latent $z$ (from either the prior or the encoder posterior). The standard REINFORCE objective is
\[
  J(\theta) \;=\; \mathbb{E}_{\tau \sim \pi_\theta}\!\left[ R(\tau) \right],
\]
with the score-function gradient estimator, using a running-mean baseline $b$ for variance reduction (as implemented):
\[
  \nabla_\theta J(\theta) \;\approx\; \mathbb{E}\!\left[ \left(R(\tau) - b\right) \cdot \nabla_\theta \log \pi_\theta(\tau) \right],
  \qquad
  \log \pi_\theta(\tau) = \sum_{t} \log \pi_\theta(a_t \mid s_t).
\]
Code correspondence: \code{logprobs} $= \sum_t \log \pi_\theta(a_t|s_t)$ (\code{logprob\_sum} in \code{new\_rl\_policy.py}), \code{advantage} $= R(\tau) - b$, and \code{rl\_loss} $= -\text{mean}[\text{advantage} \cdot \text{logprobs}] - \beta \cdot \text{mean}[H]$. This term requires the score-function trick because $R(\tau)$ is \emph{not} differentiable w.r.t.\ $\theta$ --- it passes through discrete token sampling and an external, non-differentiable predictor.

\subsection{Entropy-regularized objective}

With an entropy bonus of weight $\beta$ (\code{entropy\_coef}), the objective becomes the standard maximum-entropy RL form:
\[
  J_H(\theta) \;=\; \mathbb{E}[R(\tau)] \;+\; \beta \cdot \mathbb{E}\!\left[ H\!\left(\pi_\theta(\cdot \mid s_t)\right) \right],
\]
where $H(p) = -\sum_a p_a \log p_a$ is the per-step categorical entropy over the vocabulary, averaged across the sequence (\code{entropy\_mean} in the code). Unlike the reward term, $H$ is computed in closed form directly from the softmax output (\code{torch.distributions.Categorical(probs).entropy()}), so its gradient flows directly through the network --- no score-function estimator, no sampling noise.

\subsection{Why entropy regularization matters, in principle}

\begin{itemize}
  \item Vanilla REINFORCE ($\beta=0$) has no term that explicitly rewards diversity. Its gradient increases $\log \pi_\theta(\tau)$ for any trajectory with above-average reward and decreases it for below-average ones --- a ``rich get richer'' dynamic with no counter-force. Generically, the nearest local optimum of this dynamic is a near-deterministic policy: once some sequence scores above the running baseline, repeated reinforcement drives its probability toward $1$ and everything else toward $0$.
  \item The fixed, asymmetric \code{invalid\_penalty} amplifies this further: any sequence already verified valid is reinforced sharply relative to any newly-explored (and possibly invalid) alternative, accelerating collapse onto whatever small ``safe set'' the policy discovers first.
  \item The entropy bonus ($\beta H$) is the standard counter-force in this framework (Williams, 1992; Mnih et al., 2016; Haarnoja et al., 2018): it directly rewards the network for keeping each step's distribution closer to uniform, opposing the reward term's pressure to sharpen.
\end{itemize}

\subsection{Why raising $\beta$ after collapse had already started did not help --- a saturation argument}
\label{sec:8.4}

For a per-step categorical distribution $p = \mathrm{softmax}(z)$ over logits $z = (z_1, \ldots, z_K)$, we derive $\partial H/\partial z_j$ from first principles. Starting from $H(p) = -\sum_a p_a \log p_a$:
\begin{align}
  \frac{\partial H}{\partial z_j}
  &= -\sum_a \left[ \frac{\partial p_a}{\partial z_j} \log p_a \;+\; p_a \cdot \frac{1}{p_a}\frac{\partial p_a}{\partial z_j} \right]
  = -\sum_a \frac{\partial p_a}{\partial z_j}\left(\log p_a + 1\right). \label{eq:step1}
\end{align}
Using the softmax Jacobian $\dfrac{\partial p_a}{\partial z_j} = p_a\left(\delta_{aj} - p_j\right)$ and substituting into \eqref{eq:step1}:
\begin{align}
  \frac{\partial H}{\partial z_j}
  &= -\sum_a p_a(\delta_{aj} - p_j)(\log p_a + 1) \notag\\
  &= -p_j(\log p_j + 1) \;+\; p_j \sum_a p_a(\log p_a + 1) \notag\\
  &= -p_j(\log p_j + 1) \;+\; p_j\left[-H + 1\right] \notag\\
  &= -p_j \log p_j \;-\; p_j \;+\; p_j \;-\; p_j H \notag\\[2pt]
  &= \boxed{\;-\,p_j\bigl(\log p_j + H\bigr)\;} \label{eq:entropygrad}
\end{align}
where we used $\sum_a p_a = 1$ and $\sum_a p_a \log p_a = -H$.

As the distribution saturates toward a one-hot point (some $p_{j^*} \to 1$, all others $\to 0$, so $H \to 0$), Equation~\eqref{eq:entropygrad} gives:
\begin{itemize}
  \item For the dominant token $j^*$: $p_{j^*} \to 1$, $\log p_{j^*} \to 0$, $H \to 0$ $\;\Longrightarrow\;$ $\partial H/\partial z_{j^*} \to -1\cdot(0+0) = 0$.
  \item For every suppressed token $j$: $p_j \to 0$ dominates the divergent $\log p_j \to -\infty$ term (since $p\log p \to 0$ as $p\to 0^+$), and $H\to 0$ $\;\Longrightarrow\;$ $\partial H/\partial z_j \to 0$.
\end{itemize}

\textbf{Conclusion:} the gradient of the entropy term vanishes exactly in the regime (near-deterministic, low-$H$ policy) where it is most needed --- structurally analogous to vanishing gradients in a saturated sigmoid/softmax. Scaling $\beta$ up after this point scales an already-vanishing gradient; it cannot restore diversity once the per-step distribution has saturated.

This gives a rigorous explanation for the empirical v1$\to$v2 observation: \code{entropy\_coef} was raised 5x ($0.01\to0.05$), yet entropy collapsed further, not less. By the time the larger $\beta$ took effect, the per-step softmax outputs were already entering the saturated regime where $\partial H/\partial z \approx 0$.

\subsection{Formal implication for training design}
\label{sec:8.5}

\begin{itemize}
  \item Entropy regularization must act \emph{before} the per-step distribution saturates to have a restorative effect. Raising $\beta$ late in training, once $H$ is already small, has diminishing-to-zero returns, because both $H$ and $\partial H/\partial z$ vanish together --- a structural property of the softmax entropy gradient, not a tuning failure specific to this run.
  \item A robust fix requires a term whose gradient does not vanish under saturation --- e.g., an explicit batch-level diversity/uniqueness penalty (computed on aggregate output statistics across a batch, not the local softmax Jacobian), or scheduling a much larger $\beta$ from the very start of RL fine-tuning, or periodically raising the decode temperature to reinject exploration.
\end{itemize}

\subsection{Proposed fix: batch-level repetition penalty}
\label{sec:penalty}

A more robust alternative folds uniqueness directly into the reward signal itself, at the batch level, rather than relying on a per-step distributional statistic. For a batch of $N$ sampled sequences $\tau_1, \ldots, \tau_N$ with predictor rewards $R_1, \ldots, R_N$:
\begin{align}
  c_i &= \begin{cases} \mathrm{canonical}(\tau_i) & \text{if } \tau_i \text{ is valid} \\ \bot & \text{(invalid marker) otherwise} \end{cases} \\
  \mathrm{count}_i &= \left| \{\, j : c_j = c_i,\ c_i \neq \bot \,\} \right| \quad \text{(how many batch members share $i$'s exact molecule)} \\
  R_i^{\mathrm{div}} &= \begin{cases}
    \text{invalid\_penalty} & \text{if } \tau_i \text{ invalid} \\
    R_i - \lambda_{\mathrm{div}} \cdot \max(0,\ \mathrm{count}_i - 1) & \text{if } \tau_i \text{ valid}
  \end{cases}
\end{align}
The first occurrence of any molecule in the batch is unpenalized ($\mathrm{count}_i = 1 \Rightarrow$ no penalty); every additional duplicate accumulates penalty. Validity is still governed entirely by \code{invalid\_penalty}; this term only acts on the uniqueness axis, among already-valid samples.

\textbf{Caveat (for rigor):} this does not fully escape Section~\ref{sec:8.4}'s saturation problem --- it still reaches $\theta$ through the same score-function term $\nabla_\theta \log \pi_\theta(\tau)$, whose own Jacobian ($\partial \log p_a/\partial z_j = \delta_{aj} - p_j$) also flattens as the policy saturates. The practical advantage is \emph{magnitude and timing, not immunity}: $\mathrm{count}_i$ can be large (dozens of duplicates), giving this term far more weight per step than the bounded entropy bonus.

\subsection{Empirical test (v3) and a revised mechanism: magnitude mismatch, not only saturation}
\label{sec:8.7}

Section~\ref{sec:8.5} predicted that applying a much larger \code{entropy\_coef} from the very start of training (before any saturation) should preserve diversity. v3 tests this directly (Section~\ref{sec:v3traj}): $\beta = 0.5$ (50x v1's $0.01$), applied from epoch 1.

\textbf{Result: the prediction was falsified.} Entropy was already down to $0.10$--$0.12$ by epoch 500--1000 --- the same range v1 reached only by its final epoch $\sim$1000. Early application alone does not prevent the collapse.

The missing piece is \emph{magnitude}, not just timing. The entropy bonus is bounded above by $\beta \ln|V|$ (maximum entropy of a uniform categorical over the vocabulary, $|V|=52$ tokens here):
\[
  \max\bigl(\beta H\bigr) \;=\; \beta \ln|V| \;=\; 0.5 \times \ln(52) \;\approx\; 0.5 \times 3.95 \;\approx\; 1.98.
\]
The REINFORCE reward term has no such ceiling --- it scales with sequence length and the advantage's magnitude. The actual v3 training log shows this gap directly:

\begin{center}
\begin{tabular}{@{}r r r r@{}}
\toprule
\textbf{Epoch} & \textbf{$|$rl\_loss$|$} & \textbf{max possible $\beta H$ ($\beta=0.5$)} & \textbf{Ratio} \\
\midrule
1 & 993.5 & 1.98 & $\sim$500:1 \\
100 & 11.9 & 1.98 & $\sim$6:1 \\
500 & 1.69 & 1.98 & $\sim$1:1 \\
2000 & 2.27 & 1.98 & $\sim$1:1 \\
\bottomrule
\end{tabular}
\end{center}

\textbf{Revised mechanism:} at epoch 1, the reward term outweighs even the theoretical maximum entropy bonus by $\sim$500x --- this drives the fast initial sharpening (entropy $1.66\to0.62\to0.12$ within the first 500 epochs), and no value of $\beta$ within a reasonable range can compete at this stage. By the time $|\text{rl\_loss}|$ shrinks to a comparable scale ($\text{epoch}\sim500{+}$), the entropy bonus should matter more in relative terms --- but by then the per-step distribution has already saturated (Section~\ref{sec:8.4}), so $\partial H/\partial z$ has shrunk in lockstep. \emph{The two mechanisms compound}: magnitude mismatch causes the initial collapse (epoch 0--500), and gradient saturation then locks it in, closing the window before the entropy term's relative weight could ever help.

\textbf{Implication:} this strengthens the case for the repetition penalty (Section~\ref{sec:penalty}) over a larger \code{entropy\_coef}. Unlike $\beta H$, $R^{\mathrm{div}}$'s penalty term ($\lambda_{\mathrm{div}} \cdot \mathrm{count}_i$) is \emph{not} bounded by $\ln|V|$ --- $\mathrm{count}_i$ can scale into the dozens for a badly-collapsed batch, giving it a chance to actually compete with the reward term's magnitude at exactly the point where it is needed.

\section{Baseline (Pre-RL) Diversity Check}
\label{sec:baseline}

Evaluates the warm-start VAE checkpoint (before any RL fine-tuning) with the same script/metrics used for v1/v2/v3, to establish whether the collapse originated in RL training or was already present in the underlying VAE.

\begin{center}
\small
\begin{tabular}{@{}l l l l l@{}}
\toprule
\textbf{Mode} & \textbf{Metric} & \textbf{Pre-RL baseline} & \textbf{v1 (post-RL)} & \textbf{v2 (post-RL)} \\
\midrule
prior & Reward-model score & $-3.30\sim-4.68$ & $-3.71\sim-4.31$ & $-3.47\sim-3.55$ \\
 & Validity & 12.8--13.8\% & 48.8--57.2\% & 58.8--65.4\% \\
 & Uniqueness & 27.7--34.8\% & 37.8--38.9\% & 24.2--27.2\% \\
posterior & Reward-model score & $-3.54\sim-4.51$ & $-3.70\sim-3.90$ & $-3.14\sim-3.24$ \\
 & Validity & 7.2--10.4\% & 97.8--100\% & 97.0--98.4\% \\
 & Uniqueness & \textbf{45.8--53.9\%} & 6.3--7.8\% & 5.1--7.0\% \\
\bottomrule
\end{tabular}
\end{center}

\textbf{Confirms the causal direction directly:} the collapse is caused by RL fine-tuning, not inherent to the VAE. Before any RL, posterior-mode uniqueness is 45.8--53.9\% --- the encoder was already producing meaningfully diverse output. After RL (v1/v2), it drops to 6.3--7.8\%, a roughly 7--8x collapse. The trade for this collapse is validity: pre-RL posterior validity is only 7.2--10.4\%, vs.\ 97--100\% post-RL.

\part{v5 --- Fixing the v4 Scaffold-Repetition Loophole}

\section{Corrected Data-Scale Analysis}
\label{sec:datascale}

An earlier back-of-envelope estimate ($51$ cell lines $\times \approx 50$ compounds $\approx 2{,}550$ total (cell line, compound) pairs) was checked against the actual training data and found significantly off:

\begin{center}
\begin{tabular}{@{}l l l@{}}
\toprule
 & \textbf{Earlier estimate} & \textbf{Actual (BREAST-136344-56786-51.txt)} \\
\midrule
Total records & $\approx 2{,}550$ & $136{,}344$ \\
Compounds / cell line & $\approx 50$ (assumed uniform) & min $21$, max $43{,}061$, mean $\approx 2{,}673$ (highly skewed) \\
\bottomrule
\end{tabular}
\end{center}

\textbf{Implication:} ``full coverage'' of every (cell line, compound) pair is not achievable at any practical epoch count --- even $12{,}000$ epochs gives
\[
  \underbrace{\frac{12{,}000}{51}}_{\approx\,235\text{ visits/cell line}} \times \underbrace{64}_{\text{draws/visit}} \;\approx\; 15{,}040 \text{ max distinct draws},
\]
which would only reach $\sim$35\% of the largest single pool (ACH-000019, $43{,}061$ compounds), assuming zero redundant draws. The realistic goal for v5 is \emph{meaningfully better} coverage, not complete coverage --- addressed by raising epochs ($2{,}000 \to 10{,}000$) and cold sampling (Section~\ref{sec:cold}), not by either change alone.

\note{Cell-line-level rotation (\code{epoch \% 51}) already visits every cell line equally regardless of its pool size, in v1--v4 as well as v5 --- the imbalance problem is specifically about which compounds get sampled WITHIN a cell line's pool, not which cell lines get trained on.}

\section{v5 Change 1: Tanimoto-Threshold Repetition Penalty}
\label{sec:tanimoto}

Replaces v4's exact canonical-SMILES-equality check with fingerprint-similarity clustering: two valid molecules in the same batch are grouped together (via union-find over all pairs) if their Morgan fingerprint (radius $=2$, $2048$ bits) Tanimoto similarity is $\geq$ \code{tanimoto\_threshold}. The penalty formula is otherwise unchanged from v4:
\[
  R_i^{\mathrm{div}} \;=\; R_i - \lambda_{\mathrm{div}} \left( |C_i| - 1 \right),
\]
where $C_i$ is the similarity cluster containing sample $i$ and $|C_i|$ its size --- the first member of a cluster is unpenalized. Union-find means clustering is \emph{transitive}: molecules $A$ and $C$ can end up in the same cluster via $A$--$B$ and $B$--$C$ both clearing the threshold, even if $\mathrm{sim}(A,C)$ itself does not.

\subsection{Calibration against the actual v4 scaffold family}

Before picking a threshold, the 7 molecules that actually dominated v4's Collapse-check output (Section~\ref{sec:scaffold}) were fingerprinted and their real pairwise Tanimoto similarities computed directly, rather than guessing a value:

\begin{center}
\begin{tabular}{@{}l r@{}}
\toprule
\textbf{Pair} & \textbf{Tanimoto similarity} \\
\midrule
\code{CCCOC(O)OC} vs \code{COC(O)OCCO} & $0.714$ \\
\code{CCCOC(O)OC} vs \code{CCCCOC(O)OC} & $0.727$ \\
\code{CCCOC(O)OCO} vs \code{CCCOC(O)OCCO} & $0.762$ \\
\code{CCCOC(O)OCCO} vs \code{OCCOC(O)OCCO} & $0.778$ \\
\code{OCCOC(O)OCCO} vs \code{OCCOC(O)OCO} & $0.824$ (max) \\
\code{OCCOC(O)OCO} vs \code{CCCCOC(O)OC} & $0.321$ (min) \\
(all 21 pairs) & range $0.32$--$0.82$ \\
\bottomrule
\end{tabular}
\end{center}

\textbf{A first-guess default of $0.85$ was tested and found to be too strict}: no pair in the real family reaches $0.85$, so that default would not have caught this exact pattern at all --- the whole point of this change.

Using union-find's transitivity, the minimum threshold at which all 7 molecules chain together into one cluster was computed directly (Kruskal-style: sort pairs descending by similarity, union until fully connected):

\begin{center}
\begin{tabular}{@{}l l@{}}
\toprule
\textbf{Merge order (similarity, descending)} & \textbf{Pair} \\
\midrule
$0.824$ & \code{OCCOC(O)OCCO} $\leftrightarrow$ \code{OCCOC(O)OCO} \\
$0.778$ & \code{CCCOC(O)OCCO} $\leftrightarrow$ \code{OCCOC(O)OCCO} \\
$0.762$ & \code{CCCOC(O)OCO} $\leftrightarrow$ \code{CCCOC(O)OCCO} \\
$0.727$ & \code{CCCOC(O)OC} $\leftrightarrow$ \code{CCCCOC(O)OC} \\
$0.714$ & \code{CCCOC(O)OC} $\leftrightarrow$ \code{COC(O)OCCO} \\
$\mathbf{0.684}$ (critical --- all 7 connected after this merge) & \code{COC(O)OCCO} $\leftrightarrow$ \code{OCCOC(O)OCCO} \\
\bottomrule
\end{tabular}
\end{center}

\textbf{Chosen default: \code{--tanimoto\_threshold 0.65}} (a small margin below the critical $0.684$). Verified directly: re-running the penalty function on the real 7-molecule family plus one unrelated molecule (phenol) at threshold $=0.65$ produces exactly the expected result --- $\text{batch\_uniqueness} = 2/8 = 0.25$ (the family as one cluster, phenol separate), each family member penalized identically (cluster size $7$), phenol untouched.

\section{v5 Change 2: Cold (Least-Drawn) Compound Sampling}
\label{sec:cold}

\code{get\_seed\_batch()} previously (v1--v4) drew uniformly at random from each cell line's compound pool every epoch. Given Section~\ref{sec:datascale}'s pool-size skew, pure uniform sampling leaves most of a large pool essentially unseen even over thousands of epochs. v5 tracks a persistent per-compound draw counter and weights sampling probability as:
\[
  \boxed{\; w_j \;=\; \dfrac{1}{\left(1 + n_j\right)^{p}} \;}
\]
where $n_j$ is compound $j$'s draw count so far and $p=$~\code{cold\_power}. Under-sampled (``cold'') compounds get preferentially drawn; once drawn, $n_j$ rises and $w_j$ drops, so the mechanism is self-correcting rather than a one-time reshuffle. $p=0$ recovers v1--v4's uniform random behavior; default is $p=1.0$.

\textbf{Verified directly:} a toy pool of $10$ items drawn $20$ times at batch size $3$ ($60$ total draws, uniform expectation $=6$ each) produced counts of $[6,6,7,7,7,6,5,7,3,6]$ --- noticeably more even than pure random sampling would produce, and the underweighted item (count $3$) would be preferentially drawn next round.

\section{v5 Run Configuration}

\begin{center}
\begin{tabular}{@{}l l p{0.4\textwidth}@{}}
\toprule
\textbf{Argument} & \textbf{Value} & \textbf{Note} \\
\midrule
Script & \code{\small new\_rl\_finetune\_v5.py} & new file, does not overwrite v4's \code{\small new\_rl\_finetune.py} \\
\code{--epochs} & $10000$ & up from $2000$ (v1--v4) \\
\code{--batch} & $64$ & unchanged \\
\code{--entropy\_coef} & $0.05$ & v2/v4's value \\
\code{--invalid\_penalty} & $-5.0$ & unchanged from v2--v4 \\
\code{--lambda\_div} & $0.1$ & unchanged from v4 \\
\code{--tanimoto\_threshold} & $0.65$ & NEW --- calibrated in Section~\ref{sec:tanimoto} \\
\code{--cold\_power} & $1.0$ & NEW --- Section~\ref{sec:cold} \\
\code{--eval\_every} & $100$ & up from $5$ \\
\code{--out\_dir} & \code{checkpoints\_gen\_rl\_full\_vae\_v5} & \\
\bottomrule
\end{tabular}
\end{center}

\section{v5 Results}

\subsection{Training-log trajectory}

Job 878954, full $10{,}000$ epochs completed within a 2-hour walltime.

\begin{center}
\begin{tabular}{@{}r r r r r@{}}
\toprule
\textbf{Epoch} & \textbf{Entropy} & \textbf{Validity} & \textbf{batch\_uniqueness (Tanimoto)} & \textbf{mean\_reward\_valid} \\
\midrule
1     & 1.993 & 17.2\% & 45.5\% & $-4.54$ \\
100   & 0.408 & 45.3\% & 82.8\% & $-4.19$ \\
500   & 0.219 & 81.2\% & 76.9\% & $-4.06$ \\
1000  & 0.234 & 85.9\% & 63.6\% & $-3.76$ \\
2000  & 0.222 & 96.9\% & 54.8\% & $-3.62$ \\
3000  & 0.192 & 98.4\% & 42.9\% & $-3.57$ \\
5000  & 0.086 & 93.8\% & 36.7\% & $-3.14$ \\
7000  & 0.153 & 95.3\% & 34.4\% & $-3.39$ \\
9000  & 0.068 & 98.4\% & 36.5\% & $-3.32$ \\
10000 (tail avg, last 100) & 0.137 & 97.4\% & 35.7\% & $-3.34$ \\
\bottomrule
\end{tabular}
\end{center}

Entropy again settled into the same collapsed range as v1--v4 (tail avg $0.137$) --- consistent with Section~\ref{sec:theory}'s finding that per-step entropy is decoupled from batch-level output diversity; validity ($97.4\%$) and reward ($-3.34$) both stayed competitive with v2/v4's best results.

\subsection{Independent evaluation}
\label{sec:v5eval}

500 samples/cell line, same 3 cell lines as v1--v4. The eval script was extended with a Tanimoto-cluster-based uniqueness metric (same union-find definition as the training-time penalty) alongside the original exact-string uniqueness:

\begin{center}
\begin{tabular}{@{}l l l@{}}
\toprule
\textbf{Mode} & \textbf{Metric} & \textbf{v5 (3-cell-line range)} \\
\midrule
prior & Reward-model score & $-3.38 \sim -3.55$ \\
      & Validity & $95.4$--$97.0\%$ \\
      & Uniqueness (exact-string) & $17.9$--$20.1\%$ \\
      & Uniqueness (Tanimoto ratio, $t=0.65$) & $12.5$--$13.4\%$ \\
      & $n$ Tanimoto clusters (of 500) & $60, 61, 65$ \\
posterior & Reward-model score & $-3.32 \sim -3.53$ \\
      & Validity & $97.2$--$98.4\%$ \\
      & Uniqueness (exact-string) & $14.0$--$18.1\%$ \\
      & Uniqueness (Tanimoto ratio, $t=0.65$) & $8.5$--$13.0\%$ \\
      & $n$ Tanimoto clusters (of 500) & $60, 63, 42$ \\
both  & Novelty rate / Lipinski pass rate & $100\%$ / $100\%$ \\
\bottomrule
\end{tabular}
\end{center}

\textbf{Caution: the Tanimoto-ratio number is NOT directly comparable across different sample sizes.} Training logged \code{batch\_uniqueness} (same $t=0.65$ definition) averaging $30$--$38\%$ on $64$-sample batches; here, on $500$-sample batches, it reads $8.5$--$13.4\%$. This is a type/token-ratio artifact, not a real regression: if the policy can reach roughly a fixed number of distinct clusters $K$, then
\[
  \text{uniqueness ratio} \;=\; \frac{K}{n_{\text{valid}}}
\]
mechanically shrinks as $n_{\text{valid}}$ grows, since $K$ grows much more slowly than $n_{\text{valid}}$. The absolute cluster \emph{count} ($42$--$65$ out of $500$) is the more meaningful number at this sample size, and even that undersells the real story --- see Section~\ref{sec:headtohead}.

\subsection{v4 vs.\ v5 head-to-head: the metric that actually matters is concentration, not count}
\label{sec:headtohead}

v4's original eval predates the Tanimoto-cluster metric. Re-running v4's checkpoint through the updated eval script gives a true apples-to-apples comparison on identical cell lines/seed/sample size:

\begin{center}
\begin{tabular}{@{}l l l l@{}}
\toprule
\textbf{Mode} & \textbf{Metric} & \textbf{v4} & \textbf{v5} \\
\midrule
prior & $n$ Tanimoto clusters (of $\sim$330) & $96, 108, 116$ & $60, 61, 65$ \\
      & Largest cluster's share of valid samples & $52.6$--$60.8\%$ ($\approx57\%$) & $14.8$--$24.3\%$ ($\approx18\%$) \\
posterior & $n$ Tanimoto clusters (of $\sim$485) & $46, 46, 52$ & $42, 60, 63$ \\
      & Largest cluster's share of valid samples & $84.8$--$86.9\%$ ($\approx86\%$) & $12.1$--$20.3\%$ ($\approx16\%$) \\
\bottomrule
\end{tabular}
\end{center}

\textbf{v4's posterior-mode output was $\approx$86\% concentrated in a single Tanimoto cluster.} E.g.\ ACH-000668: $423$ of $487$ valid samples ($86.9\%$) fell into ONE cluster. v4's exact-string uniqueness looked reasonable ($28$--$32\%$) only because that mega-cluster's members were mostly distinct \emph{strings} (chain-length/branching variants of one scaffold) padded out by a long tail of one-off ``noise'' molecules that inflate the raw cluster count without adding real accessible diversity --- v4's higher raw cluster count ($46$--$116$) than v5's ($42$--$65$) is exactly this: many singleton clusters in a thin tail, not genuine breadth.

\textbf{v5's largest cluster never exceeds $\approx$24\% of any batch.} No single template dominates the way v4's scaffold family did --- generation is genuinely spread across $42$--$65$ structurally distinct clusters with no single one taking over.

\section{Final Verdict (v1--v5)}

Synthesizing Section~\ref{sec:verdict1} with the v5 results above:

\begin{itemize}
  \item \textbf{No classic reward hacking against the frozen predictor was found in any of the five runs.} Lipinski pass rate stayed $\geq$97.8\% and reward-model score stayed solidly negative throughout --- the predictor was never fooled into favoring non-drug-like molecules.
  \item \textbf{A real, controlled instance of specification gaming was found, diagnosed, and fixed} --- not against the predictor, but against the auxiliary diversity reward. v4's exact-match repetition penalty was gamed via chain-length/branching variants of one scaffold (Section~\ref{sec:scaffold}); this was traced to a concrete mechanism (exact-string matching cannot see chemical similarity) and fixed with a Tanimoto-similarity-clustered version (Section~\ref{sec:tanimoto}), calibrated against the actual failure case rather than guessed.
  \item \textbf{The metric that initially seemed to show v5 was \emph{worse} than v4 (exact-string uniqueness, 14--20\% vs.\ 28--57\%) was reversed by measuring the right thing.} Cluster \emph{concentration} --- how much of the output space one template can monopolize --- shows v5 (\pmb{$\approx$16--18\%} average) is dramatically better than v4 (\pmb{$\approx$57--86\%}). v4's higher exact-string uniqueness was itself an artifact of the very failure mode it was meant to hide.
  \item \textbf{v5 is the best result across all five runs} on the combination of validity ($\sim$97\%), reward (competitive with v2/v4's best, $-3.3\sim-3.5$), and genuine (not just nominal) diversity.
  \item \textbf{Mathematically, two independent mechanisms explain why entropy regularization alone (v1$\to$v2$\to$v3) could never have worked}, regardless of how it was tuned: (1)~the entropy bonus's gradient provably vanishes as the policy saturates (Eq.~\ref{eq:entropygrad}), and (2)~even at its theoretical maximum, the entropy bonus is bounded by $\beta\ln|V|\approx2$, while the REINFORCE reward term is unbounded and outweighs it by up to $500{:}1$ early in training. Any effective fix needed a term --- like the batch-level repetition penalty --- whose magnitude is not capped by the vocabulary size.
\end{itemize}

\end{document}
